\documentclass[lecture]{foilx}

\usepackage{latexsym}

\title{A Sample FoilX Document}
\author{FoilX}
\date{\today}
\info{Sample deck}

\newcommand{\bs}{\char '134}

\begin{document}
\maketitle

\makesection{Getting Started}

\makefoil{A First Foil}
FoilX is a small slide class inspired by FoilTeX.  It keeps the authoring
model simple: make a title page, start sections, and add foils.

\begin{itemize}
  \item Use \texttt{\bs makesection\{...\}} for section title pages.
  \item Use \texttt{\bs makefoil\{...\}} for content foils.
  \item In \texttt{lecture} mode, page numbers are shown as
        section--page.
\end{itemize}

\makefoil{Lists}
\begin{itemize}
  \item This is the first level of an itemized list.
    \begin{itemize}
      \item Here is the second level.
        \begin{itemize}
          \item And here is a third level.
          \item Nested lists use tighter spacing.
        \end{itemize}
      \item Another second-level item.
    \end{itemize}
  \item A second first-level item.
  \item A third first-level item.
\end{itemize}

\makefoil{Quotes And Enumeration}
This foil has no special page controls beyond the standard FoilX heading.

\begin{quote}
\ldots it is a good idea to make your input file easy to read.
\end{quote}

\begin{enumerate}
  \item This is an enumerated item.
  \item This is another enumerated item.
\end{enumerate}

\begin{center}
\framebox{\parbox{5.5in}{Use modern commands such as
\texttt{\bs emph\{new foil\}} rather than old font switches.}}
\end{center}

\makesection{Text And Symbols}

\makefoil{Fonts}
\begin{itemize}
  \item This shows \textrm{\textit{roman italics}}, \textsl{slanted},
        \textbf{boldface}, \texttt{typewriter}, \textrm{roman},
        and \textsc{small caps}.
  \item Unslanted \emph{emphasis} and
        \textsl{slanted \emph{emphasis}}.
  \item Size changes: {\tiny tiny}, {\scriptsize scriptsize},
        {\footnotesize footnotesize}, {\small small},
        {\normalsize normalsize}, {\large large}, {\Large Large},
        {\LARGE LARGE}, {\huge huge}, and {\Huge Huge}.
\end{itemize}

\makefoil{Special Characters And Accents}
\begin{itemize}
  \item Accents: \`{o}, \'{o}, \^{o}, \"{o}, \~{o}, \={o}, \.{o},
        \u{o}, \v{o}, \H{o}, \t{oo}, \c{o}, \d{o}, \b{o}.
  \item Paragraph symbols: \dag, \ddag, \S, \P, \pounds,
        0123456789, \copyright.
  \item Math symbols: $\dag$, $\ddag$, $\S$, $\P$, $\pounds$,
        $0123456789$.
  \item Language symbols: \oe, \OE, \ae, \AE, \aa, \AA, \o, \O,
        \l, \L, \ss, ?`, !`.
  \item TeX specials: \#, \$, \%, \&, \_, \{, \}.
\end{itemize}

\makesection{Mathematics}

\makefoil{Displayed Mathematics}
\[
{\cal F}\cdots\frac{x+y}{1+\frac{y}{z+1}}
  = \sqrt{x+y}\times \sqrt[n]{2}
\]

\[
\forall\exists\flat\natural\sharp\partial\angle\Re\mho\jmath\aleph
\]

\[
\bigcap\bigcup\bigvee\bigwedge\bigodot\bigotimes\bigoplus\biguplus
\sum\prod\coprod\int\oint
\]

\[
\sum_{i=1}^n x_i = \int_0^1 f[x]\,dx.
\]

\makefoil{Arrays}
\[
\begin{array}{clcr}
a+b+c & uv    & x-y & 27 \\
a+b   & u+v   & z   & 134 \\
a     & 3u+vw & xyz & 2,978
\end{array}
\]

\[
x=\left\{ \begin{array}{ll}
             y   & \mbox{if $y>0$} \\
             z+y & \mbox{otherwise}
          \end{array}
   \right.
\]

\makefoil{Numbered Equations}
Here is a numbered equation:
\begin{equation}
  E=mc^2
\end{equation}

And a numbered array:
\begin{eqnarray}
 x & = & 17y \\
 y & > & a+b+c+d+e+f+\nonumber \\
   &   & k+l+m+n+o+p
\end{eqnarray}

\makesection{Tables And Figures}

\makefoil{Tables}
Here is a short table:

\begin{center}
\begin{tabular}{|l|c|r|} \hline
First stuff & Second stuff & Third stuff \\ \hline\hline
foo & bar & bug \\ \hline
foofoo & barbar & bugbug \\ \hline
\end{tabular}
\end{center}

\makefoil{Picture Environments}
\begin{center}
\setlength{\unitlength}{.08in}
\begin{picture}(50,50)(0,0)
\put(10,35){\framebox(10,10){gnat}}
\put(30,20){\circle{20}}
\put(30,20){\vector(0,1){2}}
\put(45,35){\circle*{5}}
\put(30,20){\oval(20,30)}
\end{picture}
\end{center}

For real graphics, use \texttt{graphicx}, which FoilX loads by default.

\makefoil{Footnotes And Citations}
A footnote can be placed in slide text.\footnote{This is a footnote.}
This foil also cites references such as \cite{rocky} and \cite{bullwinkle}.

\begin{thebibliography}{99}

\bibitem{rocky} Rocky and Bullwinkle, Open problems, in
\textsl{Mr. Know-it-all's Rock Encyclopedia}.

\bibitem{bullwinkle} Bullwinkle, Getting things out of hats,
\textsl{Annals} \textbf{1} (1990BC), 1--2.

\end{thebibliography}

\makeappendix

\makefoil{Backup Material}
Appendix foils are numbered A1, A2, and so on.

\vf

The \texttt{\bs vf} helper inserts stretchable vertical space.

\end{document}
